\documentclass[a4paper,12pt]{article}

\usepackage[italian]{babel}
\usepackage[utf8]{inputenc}
\usepackage{geometry}
\usepackage{amsmath}
\usepackage{array}
\usepackage{circuitikz}

\geometry{margin=2.5cm}

\title{Verifica di Informatica – Versione C}
\author{Prof. Fedeli Massimo}
\date{}

\begin{document}
	
	\maketitle
	
	Nome e Cognome: \rule{8cm}{0.4pt}
	
	Classe: \rule{3cm}{0.4pt} \hfill Data: \rule{3cm}{0.4pt}
	
	\vspace{1cm}
	
	\section*{Esercizio 1 – Algoritmi e Flowgorithm}
	
	Scrivi un algoritmo in \textbf{Flowgorithm} che prenda in input dall'utente un certo numero di valori interi.
	
	Il numero di valori viene indicato dall'utente ma deve essere \textbf{almeno 10}.  
	Se l'utente inserisce un valore minore di 10 l'algoritmo deve richiederlo nuovamente.
	
	I valori devono essere memorizzati all'interno di un \textbf{vettore di interi}.
	
	Una volta inseriti tutti i numeri, l'algoritmo deve determinare:
	
	\begin{itemize}
		\item il \textbf{prodotto} di tutti i numeri
		\item il \textbf{valore massimo}
		\item il \textbf{valore minimo}
		\item la \textbf{differenza tra la somma dei numeri pari e la somma dei numeri dispari}
		\item il \textbf{numero che compare più volte} nel vettore
	\end{itemize}
	
	Corredare l'algoritmo con opportuni \textbf{commenti esplicativi}.
	
	\vspace{1cm}
	
	\section*{Esercizio 2 – Funzioni logiche}
	
	Data la seguente tabella di verità con tre variabili logiche \(A\), \(B\) e \(C\):
	
	\begin{center}
		
		\begin{tabular}{|c|c|c|c|}
			\hline
			A & B & C & F \\
			\hline
			0 & 0 & 0 & 1 \\
			0 & 0 & 1 & 0 \\
			0 & 1 & 0 & 1 \\
			0 & 1 & 1 & 0 \\
			1 & 0 & 0 & 0 \\
			1 & 0 & 1 & 1 \\
			1 & 1 & 0 & 0 \\
			1 & 1 & 1 & 1 \\
			\hline
		\end{tabular}
		
	\end{center}
	
	Determinare:
	
	\begin{enumerate}
		\item la funzione logica \(F(A,B,C)\) in \textbf{forma canonica};
		\item la \textbf{semplificazione} della funzione logica;
		\item il \textbf{circuito logico equivalente} utilizzando porte AND, OR e NOT.
	\end{enumerate}
	
	\vspace{1cm}
	
	
	Tempo a disposizione: 60 minuti
	
\end{document}